\documentclass[a4paper,12pt]{article}
\usepackage{amsmath, amsthm, amssymb,  wasysym, verbatim, bbm, color, graphics, geometry, dirtytalk, 
multicol,
algorithm,
%algorithm2e,
%algorithmic,
}
%\usepackage{calrsfs}
\usepackage{algorithm}
\usepackage[noend]{algpseudocode} % suppresses "end while", etc.
\usepackage{algorithmicx}
\algrenewcommand\algorithmicindent{2.0em} %
\usepackage[spanish]{babel}
\usepackage{marginnote}
\geometry{tmargin=.75in, bmargin=.75in, lmargin=.75in, rmargin = .75in}  
%Configurar manualmente la primera sección
 \setcounter{section}{0}
\usepackage{dirtytalk}
\usepackage{tikz}
\renewcommand\thempfootnote{\arabic{mpfootnote}} 
\renewcommand{\contentsname}{Contenido}
\newcommand{\A}{\mathcal{A}}
\newcommand{\B}{\mathcal{B}} 
\newcommand{\C}{\mathcal{C}} 
\newcommand{\D}{\mathcal{D}}
\newcommand{\R}{\mathbb{R}}
\newcommand{\N}{\mathbb{N}}
\renewcommand{\P}{\mathcal{P}}
\newcommand{\dde}{\Delta}
\newcommand{\ala}{^}
\renewcommand{\t}{\text}
\renewcommand{\r}{\rightarrow}
% \newcommand{\s}{\subseteq}
\renewcommand{\u}{\unlhd}
\newcommand{\tr}{\vartriangleleft}
\newcommand{\re}{\restriction}
\newcommand{\pr}{\preccurlyeq}
\newcommand{\prefacename}{Prefacio}
\newtheorem{lemma}{Lema}[section]
\newtheorem{proposition}{Proposici\'on}[section]
\newtheorem{theorem}{Teorema}[section]
\newtheorem*{theorem*}{Teorema}
\newtheorem{corollary}{Corolario}[section]
\newtheorem{fact}{Hecho}[section]
\newtheorem{claim}{Afirmaci\'on}[section]

\theoremstyle{definition}
\newtheorem{example}{Ejemplo}[section]
\newtheorem{example*}{Ejemplo}
\newtheorem{definition}{Definici\'on}[section]
\newtheorem{remark}{Nota}[section]
\newtheorem*{remark*}{Nota}
\newtheorem{eje}{Ejercicio}[section]
\newtheorem{question}{Pregunta}[section]
\newtheorem{notation}{Notaci\'on}[section]
\newtheorem*{notation*}{Notaci\'on}

\newtheorem{thm}{Theorem}
\newtheorem{defn}{Definition}
\newtheorem{conv}{Convention}
\newtheorem{rem}{Remark}
\newtheorem{lem}{Lemma}
\newtheorem{cor}{Corollary}

\title{Lista de ejercicios}
\author{}
\date{}

\begin{document}

\maketitle
% \tableofcontents
\vspace{.25in}

\begin{enumerate}
    \item Muestra cómo puedes traducir la información de los siguientes objetos como cadenas binarias, de modo que ser un \textit{input} en una máquina de Turing:
        \begin{enumerate}
            \item Una gráfica binaria (no ponderada) y no dirigida.
            \item Una gráfica ponderada y no dirigida.
            \item Una gráfica binaria y dirigida.
            \item Una gráfica ponderada y dirigida.
        \end{enumerate}

    \item Muestra que las siguientes operaciones son funciones computables, dando la descripción detallada de una máquina de Turing que las compute.
        \begin{enumerate}
            \item La operación de suma entre números naturales.
            \item La operación de multiplicación entre números naturales.
            \item La operación de \textit{duplicar una cadena}, i.e. si se le aplica la operación a la cadena $w$, el resultado es la cadena $ww$.
            \item La operación de comparar dos cadenas. Esta operación recibe dos cadenas como entrada, y arroja un 1 como salida si las cadenas son iguales, y un 0 si son diferentes. 
            \item La operación que verifica si una cadena es un palíndromo.
        \end{enumerate}

        \item Para este ejercicio, primero damos la siguiente definición:
        \begin{definition}\label{oblivous tm}
        Decimos que una máquina de Turing es \textit{olvidadiza} si las transiciones de sus estados no dependen del contenido de una entrada, solamente de la longitud de la misma. Es decir, una máquina $M$ es olvidadiza si y solo si, para cualesquiera cadenas de entrada $x, y\in \{0,1\}^*$, si $\lvert x\rvert = \lvert y \rvert,$ entonces $M(x)=M(y)$.
        \end{definition}
        Muestra que para toda función construible en tiempo $T:\mathbb{N}\longrightarrow \mathbb{N}$, si el lenguaje $L$ pertenece al conjunto $\mathbf{DTIME}(T(n))$, entonces existe una máquina de Turing olvidadiza que decide $L$ en tiempo $O(T^2)$.
        \item Muestra que para cada función construible en tiempo $T:\mathbb{N}\longrightarrow \mathbb{N}$, si $L\in \mathbf{DITME}(T(n))$, entonces hay una máquina de Turing que decide $L$ en tiempo $O\big(T(n)\log T(n)\big)$.

        \item Muestra que los siguientes lenguajes pertenecen a la clase $\mathbf{P}$:
        \begin{enumerate}
            \item El lenguaje de los números pares:
            \[L:=\big\{x\in \{0,1\}^*:\exists n\in \mathbb{N}\,\big(n\textit{ es par }\wedge x=Bin(n)\big)\big\}.\]
            \item El lenguaje de los números primos \[L:=\big\{x\in \{0,1\}^*:\exists n\in \mathbb{N}\,\big(n\textit{ es primo }\wedge x=Bin(n)\big)\big\}.\]
            \item El lenguaje $\texttt{CONNECTED}\subseteq\{0,1\}^*$ definido por la siguiente regla:$x\in \texttt{CONNECTED}$ si y sólo si $x$ es la representación como cadena binaria de una gráfica no dirigida $G=(V,E)$ en la que para cualesquiera vértices $x,y\in V$, existe un camino que conecta a $x$ con $y$. 
        \end{enumerate}
        \item Muestra que para cada función construible $T:\mathbb{N}\longrightarrow \mathbb{N}$, si $L\in \mathbf{DTIME}(T(n))$, entonces existe una máquina de Turing olvidadiza que decide $L$ en tiempo $O(T(n)\log T(n))$.
\end{enumerate}
\end{document}